\documentclass[a4paper, 11pt]{article}

\usepackage[T1,T2A]{fontenc}
\usepackage[utf8]{inputenc}
\usepackage[english, russian]{babel}

\usepackage{comment}
\usepackage{fullpage}
\usepackage{amssymb}
\usepackage{amsmath}
\usepackage{amsthm}
\usepackage{enumitem}

\begin{document}

\begin{center} \Large \bf Листок 13\\ Тема 13 (4.3). Лемма Гензеля, принцип Хассе \end{center}

\noindent\textbf{Упражнения и задачи}
\begin{enumerate}[topsep=0pt]

    \item Пусть $k$ --- произвольное поле, $F(X) \in k[X]$. Докажите, что существует многочлен $G(X,Y)\in k[X,Y]$, для которого $F(X+Y)=F(X)+YF'(X)+Y^2G(X,Y)$.

    \item Докажите, что порядок факторгруппы $\mathbb{Q}_2^*/(\mathbb{Q}_2^*)^2$ равен $8$. Укажите соответствующее множество представителей.

    \item Пусть $p \neq 2$, $\alpha, \beta \in \mathbb{Z}_p^\times$. Докажите, что разрешимость сравнения $\alpha x^p \equiv \beta \pmod{p^2}$ по $x\in\mathbb{Z}_p$ достаточна для разрешимости уравнения $\alpha x^p = \beta$ в $\mathbb{Q}_p$. %БШ задача 3 стр 58

    \item Докажите $p$-адический признак Эйзенштейна: если $f(x)=a_0x^n+a_1x^{n-1}+\dots+a_n\in\mathbb{Z}_p[x]$, $n\ge1$, причём $a_0\in\mathbb{Z}_p^\times$, $p\mid a_i$ для $1\le i\le n$, $p^2\nmid a_n$, то $f$ неприводим над $\mathbb{Q}_p$.

    \item Верно ли, что многочлен $f\in\mathbb{Z}[x]$ неприводим над $\mathbb{Q}$ тогда и только тогда, когда он неприводим над $\mathbb{Q}_p$ для каждого простого $p$ и над $\mathbb{R}$? %G131

    \item Докажите, что уравнение $(x^2-2)(x^2-17)(x^2-34)=0$ разрешимо в $\mathbb{Q}_p$ для каждого простого $p$ и в $\mathbb{R}$, но не разрешимо в $\mathbb{Q}$.
    
    \item Пусть $p \neq 2$, $U=U(\mathbb{Z}_p)$ --- группа $p$-адических единиц, $U_n = 1+p^n \mathbb{Z}_p$ при $n\ge1$. ($U_n = \{\alpha \in \mathbb{Z}_p: \nu_p(\alpha - 1) \geqslant n \}$, т.е. $U_n$ --- $p$-адические окрестности единичного элемента). Докажите следующие свойства:
    \begin{itemize}[topsep=0pt]
        \item $U_n/U_{n+1} \cong \mathbb{Z}/p\mathbb{Z}$, $U\cong\varprojlim_{n} U/U_n$;
        \item $U=U_1\times V$, где $V$ --- циклическая подгруппа корней степени $p-1$ из единицы;
        \item Если $\alpha \in U_n \setminus U_{n+1}$, то $\alpha^p \in U_{n+1} \setminus U_{n+2}$;
        \item $U_1/U_n$ --- циклическая группа, $U_1 \cong (\mathbb{Z}_p,+)$.
    \end{itemize}
    Сделайте вывод о структуре мультипликативной группы $\mathbb{Q}_p$: $\mathbb{Q}_p^* \cong \mathbb{Z} \times \mathbb{Z}_p \times \mathbb{Z}/(p-1)\mathbb{Z}$.
    %https://old.mccme.ru/ium/postscript/f10/AZ-Problems_3.pdf

    \item Пусть $F(x_1,\dots,x_n) \in \mathbb{Z}_p[x_1,\dots,x_n]$. При $m\ge1$ обозначим через $N_m$ число решений сравнения $F(x_1,\dots,x_n)\equiv0\pmod{p^m}$ в $(\mathbb{Z}/p^m\mathbb{Z})^n$, положим $N_0=1$. Ряд $L_F(u)=\sum_{m=0}^\infty N_m u^m$ называется рядом Пуанкаре (вспомните дзета-функцию Артина).
    \begin{itemize}[topsep=0pt]
        \item Найдите ряд Пуанкаре для $F = \alpha_1 x_1 + \dots + \alpha_n x_n$, где $\alpha_i \in \mathbb{Z}_p^*$. Убедитесь, что в этом случае $L_F(u)$ --- рациональная функция.
        \item Найдите ряд Пуанкаре для многочлена $F(x_1,\dots, x_n)$, обладающего свойством: для всякого решения сравнения $F(x_1,\dots, x_n) \equiv 0 \pmod{p}$ при некотором $i$ имеем $F'_{x_i}(x_1,\dots, x_n) \not\equiv 0 \pmod{p}$;
        \item Найдите ряд Пуанкаре для $F(x,y)=x^2-y^3$;
        \item Докажите рациональность ряда Пуанкаре для многочленов одной переменной.
    \end{itemize}
    (Существует теорема (Игусы) о том, что $L_F$ всегда является рациональной).
    %БШ задачи 9-12 стр 59

\end{enumerate}

% \noindent\textbf{SageMath}
% \begin{itemize}[topsep=0pt]
%     \item Исследуйте функции SageMath для работы с многочленами с $p$-адическими коэффициентами:
%     \begin{itemize}[noitemsep,topsep=0pt]
%         \item Разложение многочлена: \texttt{factor()}.
%         \item В контексте Леммы Гензеля: \texttt{hensel\_lift()}.
%     \end{itemize}
% \end{itemize}

\noindent\textbf{Темы для самостоятельного изучения}
\begin{itemize}[topsep=0pt]
    \item Разрешимость уравнений с квадратичными формами над полем $p$-адических чисел ([БШ, \S I.6, п.~2]).
    \item Приложение теоремы Минковского-Хассе для квадратичной формы от трёх переменных ([Gouv, \S 4.8]).
\end{itemize}

\end{document}